\begin{definition}[Controlled dependent partial trace]
    \label{def:mpdo_controlled_dependent_partial_trace}
    \lean{Matrix.controlledPartialTraceRightLM}
    \leanok
    \uses{def:mpdo_controlled_kraus_map, def:mpdo_right_partial_trace_map}
    Let
    \begin{align}
        H&=\bigoplus_{i\in I}(A_i\otimes B_i), \notag\\
        K&=\bigoplus_{i\in I}A_i. \notag
    \end{align}
    The controlled dependent partial trace discards the off-diagonal blocks
    between distinct $i$ and applies $\tr_{B_i}$ to the $i$th diagonal block.
\end{definition}

\begin{theorem}[Diagonal blocks of the controlled dependent partial trace]
    \label{thm:mpdo_controlled_dependent_partial_trace_same_block}
    \lean{Matrix.controlledPartialTraceRightLM_sameBlock_apply}
    \leanok
    \uses{def:mpdo_controlled_dependent_partial_trace,
      def:mpdo_right_partial_trace_map}
    If $X_{ii}$ denotes the $i$th diagonal block of $X$, then
    $[\mathcal C_{\tr}(X)]_{ii}=\tr_{B_i}(X_{ii})$.
\end{theorem}
\begin{proof}
    \leanok
    \uses{thm:mpdo_controlled_kraus_map_same_block}
    If $(K_{i,j})_j$ is the Kraus family chosen for $\tr_{B_i}$, then the
    diagonal-block formula gives
    \begin{align}
        [\mathcal C_{\tr}(X)]_{ii}
        &=\sum_j K_{i,j}X_{ii}K_{i,j}^\dagger
          =\tr_{B_i}(X_{ii}). \notag
    \end{align}
\end{proof}

\begin{theorem}[The controlled dependent partial trace is trace-preserving
    completely positive]
    \label{thm:mpdo_controlled_dependent_partial_trace_cptp}
    \lean{Matrix.controlledPartialTraceRightLM_isKrausCPTP}
    \leanok
    \uses{def:mpdo_controlled_dependent_partial_trace}
    The controlled dependent partial trace is trace-preserving and completely
    positive.
\end{theorem}
\begin{proof}
    \leanok
    \uses{thm:mpdo_controlled_kraus_map_cptp,
      lem:mpdo_right_partial_trace_cptp}
    Choose Kraus operators $(V_{i,a})_a$ for $\tr_{B_i}$ and extend each
    $V_{i,a}$ by zero outside the $i$th summand. Their controlled family
    satisfies
    \begin{align}
        \sum_{i,a}\widetilde V_{i,a}^{\dagger}\widetilde V_{i,a}
        &=\bigoplus_i\left(\sum_aV_{i,a}^{\dagger}V_{i,a}\right)
          =\bigoplus_i \Id_{A_i\otimes B_i}
          =\Id_H. \notag
    \end{align}
    Its Kraus form gives complete positivity, and the displayed resolution of
    the identity gives trace preservation.
\end{proof}

\begin{definition}[Two-site subspin regrouping]
    \label{def:mpdo_two_site_subspin_regrouping}
    \lean{MPOTensor.ExplicitEtaOperators.PhysicalSectorIndex}
    \lean{MPOTensor.ExplicitEtaOperators.OuterSubspinIndex}
    \lean{MPOTensor.ExplicitEtaOperators.NeighborSubspinIndex}
    \lean{MPOTensor.ExplicitEtaOperators.SectorSiteIndex}
    \lean{MPOTensor.ExplicitEtaOperators.BoundarySubspinIndex}
    \lean{MPOTensor.ExplicitEtaOperators.EtaPreparationIndex}
    \lean{MPOTensor.ExplicitEtaOperators.twoSiteTraceRegroupEquiv}
    \lean{MPOTensor.ExplicitEtaOperators.t0RegroupEquiv}
    \lean{MPOTensor.ExplicitEtaOperators.etaShiftEquiv}
    \lean{MPOTensor.ExplicitEtaOperators.s2ShiftEquiv}
    \lean{MPOTensor.PhysicalSectorFactorization.t0RegroupEquiv}
    \leanok
    \uses{def:mpdo_eta_structure, def:mpdo_physical_sector_factorization}
    In sector-adapted coordinates, regroup two sites by
    \begin{align}
        ((l_k,r_k),(l_h,r_h))
        &\longmapsto ((l_k,r_h),(r_k,l_h)). \notag
    \end{align}
    Globally this gives the dependent-sum equivalence
    \begin{align}
        \left(\bigoplus_k L_k\otimes R_k\right)^{\!\otimes2}
        &\simeq
          \bigoplus_{k,h}(L_k\otimes R_h)\otimes(R_k\otimes L_h). \notag
    \end{align}
    The forward equivalence prepares the neighboring factors
    $R_k\otimes L_h$ for the partial trace in $\mathcal T_0$; its inverse is
    the shift $\mathcal S_2$. These are the reorderings in
    \cite[Appendix~C.2, lines~1521--1522 and~1555--1559]
      {Cirac2016MPDO_arXiv}.
\end{definition}

\begin{definition}[The two-site partial-trace map $\mathcal T_0$]
    \label{def:mpdo_t0_map}
    \lean{MPOTensor.ExplicitEtaOperators.t0Map}
    \lean{MPOTensor.PhysicalSectorFactorization.t0Map}
    \leanok
    \uses{def:mpdo_two_site_subspin_regrouping,
      def:mpdo_controlled_dependent_partial_trace,
      def:mpdo_equiv_reindex_map}
    Define $\mathcal T_0$ by the two-site regrouping followed, in each pair
    $(k,h)$, by the partial trace over $R_k\otimes L_h$. Its output is indexed
    by the retained outer factors $L_k\otimes R_h$. Thus, writing $R_{\Phi_2}$
    for the regrouping reindexing,
    \begin{align}
        \mathcal T_0
        &=\CTr_{(R_k\otimes L_h)_{k,h}}\circ R_{\Phi_2}.
        \label{eq:rfp_t0_factorization}
    \end{align}
\end{definition}

\begin{theorem}[$\mathcal T_0$ is trace-preserving completely positive]
    \label{thm:mpdo_t0_map_cptp}
    \lean{MPOTensor.ExplicitEtaOperators.t0Map_isKrausCPTP}
    \lean{MPOTensor.PhysicalSectorFactorization.t0Map_isKrausCPTP}
    \leanok
    \uses{def:mpdo_t0_map}
    The map $\mathcal T_0$ is trace-preserving and completely positive.
\end{theorem}
\begin{proof}
    \leanok
    \uses{thm:mpdo_equiv_reindex_cptp,
      thm:mpdo_controlled_dependent_partial_trace_cptp,
      lem:is_kraus_cptp_comp}
    By (\ref{eq:rfp_t0_factorization}), $\mathcal T_0$ is the composition of
    the controlled dependent partial trace with reindexing along the two-site
    equivalence. The first map is trace-preserving and completely positive by
    Theorem~\ref{thm:mpdo_controlled_dependent_partial_trace_cptp}, and the
    second by Theorem~\ref{thm:mpdo_equiv_reindex_cptp}. Their composition is
    therefore trace-preserving and completely positive.
\end{proof}

\begin{definition}[The two-site shift $\mathcal S_2$]
    \label{def:mpdo_s2_map}
    \lean{MPOTensor.ExplicitEtaOperators.s2Map}
    \lean{MPOTensor.PhysicalSectorFactorization.s2ShiftEquiv,
      MPOTensor.PhysicalSectorFactorization.s2Map,
      MPOTensor.PhysicalSectorFactorization.s2Map_apply}
    \leanok
    \uses{def:mpdo_two_site_subspin_regrouping,
      def:mpdo_equiv_reindex_map}
    The map $\mathcal S_2$ is matrix reindexing along the inverse two-site
    equivalence. Thus
    \begin{align}
        ((l_k,r_h),(r_k,l_h))
        &\longmapsto ((l_k,r_k),(l_h,r_h)). \notag
    \end{align}
\end{definition}

\begin{theorem}[$\mathcal S_2$ is trace-preserving completely positive]
    \label{thm:mpdo_s2_map_cptp}
    \lean{MPOTensor.ExplicitEtaOperators.s2Map_isKrausCPTP}
    \lean{MPOTensor.PhysicalSectorFactorization.s2Map_isKrausCPTP}
    \leanok
    \uses{def:mpdo_s2_map}
    The shift $\mathcal S_2$ is trace-preserving and completely positive.
\end{theorem}
\begin{proof}
    \leanok
    \uses{thm:mpdo_equiv_reindex_cptp}
    If $\Phi_2$ is the two-site regrouping, then
    $\mathcal S_2=R_{\Phi_2^{-1}}$. It is therefore trace-preserving and
    completely positive by Theorem~\ref{thm:mpdo_equiv_reindex_cptp}.
\end{proof}

\begin{definition}[Three-site subspin regrouping]
    \label{def:mpdo_three_site_subspin_regrouping}
    \lean{MPOTensor.ExplicitEtaOperators.OmegaPreparationIndex}
    \lean{MPOTensor.ExplicitEtaOperators.threeSiteTraceRegroupEquiv}
    \lean{MPOTensor.ExplicitEtaOperators.s0RegroupEquiv}
    \lean{MPOTensor.ExplicitEtaOperators.omegaShiftEquiv}
    \lean{MPOTensor.ExplicitEtaOperators.t2ShiftEquiv}
    \lean{MPOTensor.PhysicalSectorFactorization.ThreeSiteMiddleIndex,
      MPOTensor.PhysicalSectorFactorization.S0PreparationIndex,
      MPOTensor.PhysicalSectorFactorization.s0RegroupEquiv,
      MPOTensor.PhysicalSectorFactorization.t2ShiftEquiv}
    \leanok
    \uses{def:mpdo_two_site_subspin_regrouping, def:mpdo_eta_omega,
      def:mpdo_physical_sector_factorization}
    Regroup three sites by
    \begin{align}
        ((l_k,r_k),((l_l,r_l),(l_h,r_h)))
        &\longmapsto
          ((l_k,r_h),((r_k,l_l),(r_l,l_h))). \notag
    \end{align}
    Equivalently,
    \begin{align}
        \left(\bigoplus_k L_k\otimes R_k\right)^{\!\otimes3}
        &\simeq
          \bigoplus_{k,h}(L_k\otimes R_h)\otimes
          \left[\bigoplus_l
            (R_k\otimes L_l)\otimes(R_l\otimes L_h)\right]. \notag
    \end{align}
    The forward equivalence prepares the four middle subspins for
    $\mathcal S_0$; its inverse is the shift $\mathcal T_2$ in
    \cite[Appendix~C.2, lines~1535--1540 and~1547]
      {Cirac2016MPDO_arXiv}.
\end{definition}

\begin{definition}[The three-site partial-trace map $\mathcal S_0$]
    \label{def:mpdo_s0_map}
    \lean{MPOTensor.ExplicitEtaOperators.s0Map}
    \lean{MPOTensor.PhysicalSectorFactorization.s0Map,
      MPOTensor.PhysicalSectorFactorization.s0Map_sameBlock_apply}
    \leanok
    \uses{def:mpdo_three_site_subspin_regrouping,
      def:mpdo_controlled_dependent_partial_trace,
      def:mpdo_equiv_reindex_map}
    Define $\mathcal S_0$ by the three-site regrouping followed, for every
    outer pair $(k,h)$, by the partial trace over
    \begin{align}
        \bigoplus_l(R_k\otimes L_l)\otimes(R_l\otimes L_h). \notag
    \end{align}
    Thus, writing $R_{\Phi_3}$ for the regrouping reindexing,
    \begin{align}
        \mathcal S_0
        &=\CTr_{\left(\bigoplus_l(R_k\otimes L_l)\otimes(R_l\otimes L_h)\right)_{k,h}}\circ R_{\Phi_3}.
        \label{eq:rfp_s0_factorization}
    \end{align}
\end{definition}
% The source prose at line 1547 repeats 2l; the diagram forces the final
% traced factor to be 3l. The correction is documented in
% docs/paper-gaps/cpgsv17_mpdo_s0_trace_subspin_typo.tex.

\begin{theorem}[$\mathcal S_0$ is trace-preserving completely positive]
    \label{thm:mpdo_s0_map_cptp}
    \lean{MPOTensor.ExplicitEtaOperators.s0Map_isKrausCPTP}
    \lean{MPOTensor.PhysicalSectorFactorization.s0Map_isKrausCPTP}
    \leanok
    \uses{def:mpdo_s0_map}
    The map $\mathcal S_0$ is trace-preserving and completely positive.
\end{theorem}
\begin{proof}
    \leanok
    \uses{thm:mpdo_equiv_reindex_cptp,
      thm:mpdo_controlled_dependent_partial_trace_cptp,
      lem:is_kraus_cptp_comp}
    By (\ref{eq:rfp_s0_factorization}), $\mathcal S_0$ is the composition of
    the controlled dependent partial trace with reindexing along the three-site
    equivalence. The controlled trace and equivalence reindexing are
    trace-preserving and completely positive by
    Theorems~\ref{thm:mpdo_controlled_dependent_partial_trace_cptp}
    and~\ref{thm:mpdo_equiv_reindex_cptp}, respectively. Their composition
    has the same property.
\end{proof}

\begin{definition}[The three-site shift $\mathcal T_2$]
    \label{def:mpdo_t2_map}
    \lean{MPOTensor.ExplicitEtaOperators.t2Map}
    \lean{MPOTensor.PhysicalSectorFactorization.t2Map}
    \leanok
    \uses{def:mpdo_three_site_subspin_regrouping,
      def:mpdo_equiv_reindex_map}
    The map $\mathcal T_2$ is matrix reindexing along the inverse three-site
    equivalence. On each summand it sends
    \begin{align}
        ((l_k,r_h),((r_k,l_l),(r_l,l_h)))
        &\longmapsto
          ((l_k,r_k),((l_l,r_l),(l_h,r_h))). \notag
    \end{align}
\end{definition}

\begin{theorem}[$\mathcal T_2$ is trace-preserving completely positive]
    \label{thm:mpdo_t2_map_cptp}
    \lean{MPOTensor.ExplicitEtaOperators.t2Map_isKrausCPTP}
    \lean{MPOTensor.PhysicalSectorFactorization.t2Map_isKrausCPTP}
    \leanok
    \uses{def:mpdo_t2_map}
    The shift $\mathcal T_2$ is trace-preserving and completely positive.
\end{theorem}
\begin{proof}
    \leanok
    \uses{thm:mpdo_equiv_reindex_cptp}
    If $\Phi_3$ is the three-site regrouping, then
    $\mathcal T_2=R_{\Phi_3^{-1}}$. It is therefore trace-preserving and
    completely positive by Theorem~\ref{thm:mpdo_equiv_reindex_cptp}.
\end{proof}

\begin{definition}[Fiberwise partial-trace maps]
    \label{def:mpdo_fiberwise_partial_trace_maps}
    \lean{MPOTensor.ExplicitEtaOperators.t0SectorMap}
    \lean{MPOTensor.ExplicitEtaOperators.s0SectorMap}
    \leanok
    \uses{def:mpdo_two_site_subspin_regrouping,
      def:mpdo_three_site_subspin_regrouping,
      def:mpdo_right_partial_trace_map,
      def:mpdo_equiv_reindex_map}
    For fixed sectors $(k,h)$, define the fiberwise form of $\mathcal T_0$ by
    regrouping two sites and tracing $R_k\otimes L_h$. For fixed
    $(k,l,h)$, define the fiberwise form of $\mathcal S_0$ by regrouping three
    sites and tracing
    $(R_k\otimes L_l)\otimes(R_l\otimes L_h)$.
\end{definition}

\begin{theorem}[The fiberwise partial traces are trace-preserving completely
    positive]
    \label{thm:mpdo_fiberwise_partial_trace_maps_cptp}
    \lean{MPOTensor.ExplicitEtaOperators.t0SectorMap_isKrausCPTP}
    \lean{MPOTensor.ExplicitEtaOperators.s0SectorMap_isKrausCPTP}
    \leanok
    \uses{def:mpdo_fiberwise_partial_trace_maps}
    Every fiberwise form of $\mathcal T_0$ and $\mathcal S_0$ is
    trace-preserving and completely positive.
\end{theorem}
\begin{proof}
    \leanok
    \uses{thm:mpdo_equiv_reindex_cptp,
      lem:mpdo_right_partial_trace_cptp,
      lem:is_kraus_cptp_comp}
    The fiberwise forms factor as
    \begin{align}
        \mathcal T_0^{(k,h)}
        &=\tr_{R_k\otimes L_h}\circ R_{\phi_{k,h}}, \notag\\
        \mathcal S_0^{(k,l,h)}
        &=\tr_{(R_k\otimes L_l)\otimes(R_l\otimes L_h)}
          \circ R_{\phi_{k,l,h}}, \notag
    \end{align}
    where $\phi_{k,h}$ is the two-site equivalence of
    Definition~\ref{def:mpdo_two_site_subspin_regrouping} on the $(k,h)$
    sector, $\phi_{k,l,h}$ is the three-site equivalence of
    Definition~\ref{def:mpdo_three_site_subspin_regrouping} on the $(k,l,h)$
    sector, and $R_\phi$ denotes equivalence reindexing. Both factors in each
    composition are trace-preserving and completely positive, so the
    composition is as well.
\end{proof}

\begin{definition}[Fiberwise shifts]
    \label{def:mpdo_fiberwise_shifts}
    \lean{MPOTensor.ExplicitEtaOperators.s2SectorMap}
    \lean{MPOTensor.ExplicitEtaOperators.t2SectorMap}
    \leanok
    \uses{def:mpdo_two_site_subspin_regrouping,
      def:mpdo_three_site_subspin_regrouping,
      def:mpdo_equiv_reindex_map}
    On fixed sector pairs, define $\mathcal S_2$ and $\mathcal T_2$ by the
    same inverse regroupings as the corresponding global shifts.
\end{definition}

\begin{samepage}
\begin{theorem}[The fiberwise shifts are trace-preserving completely positive]
    \label{thm:mpdo_fiberwise_shifts_cptp}
    \lean{MPOTensor.ExplicitEtaOperators.s2SectorMap_isKrausCPTP}
    \lean{MPOTensor.ExplicitEtaOperators.t2SectorMap_isKrausCPTP}
    \leanok
    \uses{def:mpdo_fiberwise_shifts}
    Every fiberwise shift is trace-preserving and completely positive.
\end{theorem}
\begin{proof}
    \leanok
    \uses{thm:mpdo_equiv_reindex_cptp}
    The shifts $\mathcal S_2^{(k,h)}$ and $\mathcal T_2^{(k,h)}$ are
    reindexings along the inverse two-site and three-site equivalences of
    Definitions~\ref{def:mpdo_two_site_subspin_regrouping}
    and~\ref{def:mpdo_three_site_subspin_regrouping}, respectively. They are
    therefore trace-preserving and completely positive by
    Theorem~\ref{thm:mpdo_equiv_reindex_cptp}.
\end{proof}
\end{samepage}

To record the domains and codomains, set
\begin{align}
    \mathcal H^{\Omega}_{k,h}
    &=\bigoplus_l
      (B_k^R\otimes B_l^L)\otimes(B_l^R\otimes B_h^L), \notag
\end{align}
and introduce the matrix algebras
\begin{align}
    \mathcal A_2
    &=\End\!\left(\left[\bigoplus_j(B_j^L\otimes B_j^R)\right]^{\!\otimes2}\right),
    &
    \mathcal A_{\partial}
    &=\End\!\left(\bigoplus_{k,h}B_k^L\otimes B_h^R\right), \notag\\
    \mathcal A_{\eta}
    &=\End\!\left(\bigoplus_{k,h}\bigl[(B_k^L\otimes B_h^R)\otimes
          (B_k^R\otimes B_h^L)\bigr]\right),
    &
    \mathcal A_{\Omega}
    &=\End\!\left(\bigoplus_{k,h}\bigl[(B_k^L\otimes B_h^R)\otimes
          \mathcal H^{\Omega}_{k,h}\bigr]\right), \notag\\
    \mathcal A_3
    &=\End\!\left(\left[\bigoplus_j(B_j^L\otimes B_j^R)\right]^{\!\otimes3}\right).
    \notag
\end{align}
Thus the global maps have types
\begin{align}
    \mathcal T_0&:\mathcal A_2\to\mathcal A_{\partial}, \notag\\
    \mathcal S_2&:\mathcal A_{\eta}\to\mathcal A_2, \notag\\
    \mathcal S_0&:\mathcal A_3\to\mathcal A_{\partial}, \notag\\
    \mathcal T_2&:\mathcal A_{\Omega}\to\mathcal A_3. \notag
\end{align}
The controlled maps $\mathcal T_0$ and $\mathcal S_0$ discard coherences
between distinct outer-sector pairs. For fixed $(k,h)$, the global map
$\mathcal S_0$ traces the entire carrier $\mathcal H^{\Omega}_{k,h}$, including
its direct sum over $l$. The shifts $\mathcal S_2$ and $\mathcal T_2$ are
global reindexings, with all displayed tensor factors present at their inputs
\cite[Appendix~C.2, lines~1521--1559]{Cirac2016MPDO_arXiv}.

For the fiberwise maps, write
\begin{align}
    \mathcal A_2^{k,h}
    &=\End\!\left((B_k^L\otimes B_k^R)\otimes(B_h^L\otimes B_h^R)\right),
    &
    \mathcal A_{\partial}^{k,h}
    &=\End(B_k^L\otimes B_h^R), \notag\\
    \mathcal A_{\eta}^{k,h}
    &=\End\!\left((B_k^L\otimes B_h^R)\otimes(B_k^R\otimes B_h^L)\right),
    &
    \mathcal A_{\Omega}^{k,h}
    &=\End\!\left((B_k^L\otimes B_h^R)\otimes\mathcal H^{\Omega}_{k,h}\right), \notag\\
    \mathcal A_3^{k,l,h}
    &=\End\!\left((B_k^L\otimes B_k^R)\otimes(B_l^L\otimes B_l^R)
          \otimes(B_h^L\otimes B_h^R)\right), \notag\\
    \mathcal A_3^{k,\bullet,h}
    &=\End\!\left(\bigoplus_l
        \bigl[(B_k^L\otimes B_k^R)\otimes(B_l^L\otimes B_l^R)
          \otimes(B_h^L\otimes B_h^R)\bigr]\right). \notag
\end{align}
The map $\mathcal S_0^{(k,l,h)}$ fixes the middle sector $l$ before taking the
partial trace; it is therefore distinct from the fixed-$(k,h)$ restriction of
the global controlled map $\mathcal S_0$. The four fiberwise types are
% Formula: T_0^(k,h), S_2^(k,h), S_0^(k,l,h), and T_2^(k,h) have the displayed
% matrix-algebra domains and codomains. Ink: every blue stroke denotes a
% trace-preserving completely positive map, and its label specifies the map.
\par\noindent\hfil
  \begin{tikzcd}[channel, column sep=16mm, row sep=6mm]
    \mathcal A_2^{k,h} \arrow[r, "\mathcal T_0^{(k,h)}"] &
      \mathcal A_{\partial}^{k,h} \\
    \mathcal A_{\eta}^{k,h} \arrow[r, "\mathcal S_2^{(k,h)}"] &
      \mathcal A_2^{k,h} \\
    \mathcal A_3^{k,l,h} \arrow[r, "\mathcal S_0^{(k,l,h)}"] &
      \mathcal A_{\partial}^{k,h} \\
    \mathcal A_{\Omega}^{k,h} \arrow[r, "\mathcal T_2^{(k,h)}"] &
      \mathcal A_3^{k,\bullet,h}
  \end{tikzcd}
\hfil\par
